% ======================================================================
% Appendix J. Stationary phase and oscillatory integrals
% PATH: src/appendices/J-stationary-phase-and-oscillatory-integrals.tex
% ======================================================================

\appendix
\chapter{Stationary phase and oscillatory integrals}
\label{app:J.stationary}

\section{Scope and conventions}
\label{sec:J.scope}

This appendix collects the oscillatory integral estimates used repeatedly in Chapters~\ref{chap:05-microlocal}--%
\ref{chap:08-normalization} and in the application blocks.  The emphasis is on \emph{uniformity} with respect to the
two-scale parameters that appear in the localized trace framework: the large spectral parameter $\lambda$ and the
window scale $\eta$ (typically $\lambda^{-\theta}\le \eta\le 1$).  We work throughout in the Fourier conventions fixed
in Appendix~\ref{app:G.fourier}.

We use the notation $A\lesssim B$ to mean $A\le C\,B$ for a constant $C$ depending only on the ambient dimension and
on finitely many seminorm bounds for the amplitudes under consideration; dependence on additional parameters is
always stated explicitly.

\subsection{Symbol and cutoff classes}
\label{subsec:J.symbols}

Let $\langle x\rangle=(1+|x|^2)^{1/2}$.  For $m\in\R$ we denote by $S^m(\R^d)$ the standard symbol class:
$a\in S^m(\R^d)$ if $a\in C^\infty(\R^d)$ and for all multiindices $\alpha$ there exists $C_\alpha$ with
\[
|\partial_x^\alpha a(x)|\le C_\alpha \langle x\rangle^{m-|\alpha|}.
\]
When amplitudes depend on external parameters $(\lambda,\eta)$ we interpret seminorm bounds uniformly in the admissible
parameter range unless otherwise stated.

For $\eta>0$ we use the rescaled cutoff $\chi_\eta(t)=\chi(t/\eta)$, with $\chi\in\mathcal{S}(\R)$ fixed, even.
We record the standard scaling identities:
\begin{equation}\label{eq:J.scaling}
\widehat{\chi_\eta}(\xi)=\eta\,\widehat{\chi}(\eta\xi),
\qquad
\|\partial_t^k \chi_\eta\|_{L^\infty}\lesssim \eta^{-k}.
\end{equation}

\section{Integration by parts: nonstationary phase}
\label{sec:J.nsp}

We begin with the standard nonstationary phase lemma in a form suited to uniform bounds.

\begin{lemma}[Nonstationary phase]
\label{lem:J.nonstationary}
Let $\phi\in C^\infty(\R^d)$ and $a\in C_c^\infty(\R^d)$.
Assume that on $\supp(a)$ one has $|\nabla\phi(x)|\ge c_0>0$.
Then for every $N\ge 0$ and every $\lambda\ge 1$,
\begin{equation}\label{eq:J.nonstat.bound}
\left|\int_{\R^d} e^{i\lambda\phi(x)} a(x)\,dx\right|
\le
C_N\,\lambda^{-N}\sum_{|\alpha|\le N}\|\partial^\alpha a\|_{L^1},
\end{equation}
where $C_N$ depends only on $N,d$ and finitely many bounds for derivatives of $\phi$ on $\supp(a)$, and on $c_0^{-1}$.
\end{lemma}

\begin{proof}
Define the differential operator
\[
L=\frac{1}{i\lambda}\,\frac{\nabla\phi\cdot\nabla}{|\nabla\phi|^2},
\qquad
L\big(e^{i\lambda\phi}\big)=e^{i\lambda\phi}.
\]
Integrating by parts $N$ times yields \eqref{eq:J.nonstat.bound}.
\end{proof}

\begin{remark}[Uniformity under parameter dependence]
\label{rem:J.uniform.nsp}
If $\phi=\phi_{\lambda,\eta}$ and $a=a_{\lambda,\eta}$ depend on external parameters and the lower bound
$|\nabla\phi_{\lambda,\eta}|\ge c_0$ holds uniformly on $\supp(a_{\lambda,\eta})$, while finitely many derivatives of
$\phi_{\lambda,\eta}$ are uniformly bounded, then \eqref{eq:J.nonstat.bound} is uniform in $(\lambda,\eta)$.
\end{remark}

\section{One-dimensional stationary phase}
\label{sec:J.1d}

We record the classical stationary phase expansion in one dimension with explicit remainder control.

\begin{theorem}[Stationary phase in one dimension]
\label{thm:J.sp1d}
Let $\phi\in C^\infty([a,b])$ and $A\in C_c^\infty((a,b))$.
Assume $\phi'$ has a unique zero $x_0\in(a,b)$ and that $\phi''(x_0)\neq 0$.
Then for $\lambda\to\infty$ one has the expansion
\begin{equation}\label{eq:J.sp1d.expansion}
\int_a^b e^{i\lambda\phi(x)}A(x)\,dx
=
e^{i\lambda\phi(x_0)}e^{i\frac{\pi}{4}\sgn\phi''(x_0)}
\left(\frac{2\pi}{\lambda|\phi''(x_0)|}\right)^{1/2}
\sum_{k=0}^{N-1}\lambda^{-k}\,c_k
+
R_N(\lambda),
\end{equation}
where $c_k$ are universal differential expressions in $\phi$ and $A$ evaluated at $x_0$, and
\begin{equation}\label{eq:J.sp1d.remainder}
|R_N(\lambda)|
\le
C_N\,\lambda^{-N-\frac12},
\end{equation}
with $C_N$ depending on finitely many derivatives of $\phi$ and $A$ on $[a,b]$ and on $|\phi''(x_0)|^{-1}$.
\end{theorem}

\begin{remark}[Leading coefficient]
\label{rem:J.sp1d.leading}
The leading coefficient is $c_0=A(x_0)$.  In particular,
\[
\int e^{i\lambda\phi(x)}A(x)\,dx
=
e^{i\lambda\phi(x_0)}e^{i\frac{\pi}{4}\sgn\phi''(x_0)}
\left(\frac{2\pi}{\lambda|\phi''(x_0)|}\right)^{1/2}
A(x_0)
+
O(\lambda^{-3/2}).
\]
\end{remark}

\section{Multi-dimensional stationary phase}
\label{sec:J.md}

We now record the $d$-dimensional stationary phase theorem in the nondegenerate case.

\begin{theorem}[Nondegenerate stationary phase]
\label{thm:J.spd}
Let $\phi\in C^\infty(\R^d)$ and $a\in C_c^\infty(\R^d)$.
Assume $\nabla\phi$ has a unique critical point $x_0\in\supp(a)$ and that the Hessian
$H=\nabla^2\phi(x_0)$ is nondegenerate.
Then for $\lambda\to\infty$,
\begin{equation}\label{eq:J.spd.expansion}
\int_{\R^d} e^{i\lambda\phi(x)}a(x)\,dx
=
e^{i\lambda\phi(x_0)}e^{i\frac{\pi}{4}\sgn(H)}
\left(\frac{2\pi}{\lambda}\right)^{d/2}
\frac{1}{|\det H|^{1/2}}
\sum_{k=0}^{N-1}\lambda^{-k}\,C_k
+
R_N(\lambda),
\end{equation}
where $C_k$ are universal differential expressions in $\phi$ and $a$ at $x_0$, and
\begin{equation}\label{eq:J.spd.remainder}
|R_N(\lambda)|
\le
C_N\,\lambda^{-N-\frac{d}{2}},
\end{equation}
with $C_N$ depending on finitely many derivatives of $\phi,a$ on $\supp(a)$ and on a lower bound for
$|\det H|$.
\end{theorem}

\begin{remark}[Uniformity under rescaling]
\label{rem:J.spd.uniform}
In the localized trace applications, the phase typically has the form
$\phi(x)=\phi_0(x)+O(\eta)$ and $\lambda$ is large while $\eta$ is small.
Uniformity of \eqref{eq:J.spd.expansion} requires that the nondegeneracy bounds for the Hessian persist uniformly in
the admissible parameter range.  This is ensured in our applications by the microlocal localization away from caustics
and conjugate points, as in Chapter~\ref{chap:05-microlocal}.
\end{remark}

\section{Parameterized stationary phase at two scales}
\label{sec:J.two.scale}

We record a form of stationary phase suited to integrals where both $\lambda$ and $\eta$ appear, typically through
amplitudes supported on $|t|\lesssim \eta^{-1}$.

\subsection{A model estimate}
\label{subsec:J.model}

Let $I(\lambda,\eta)$ be of the form
\begin{equation}\label{eq:J.model.integral}
I(\lambda,\eta)
=
\int_{\R^d} e^{i\lambda\phi(x)} a(x)\,\chi\!\left(\frac{\psi(x)}{\eta}\right)\,dx,
\qquad
\lambda\ge 1,\quad 0<\eta\le 1,
\end{equation}
where $\chi\in C_c^\infty(\R)$ and $\psi$ is smooth with $\psi(x_0)=0$ at a stationary point $x_0$ of $\phi$.

The factor $\chi(\psi/\eta)$ forces the integration region to a tube $\{|\psi|\lesssim \eta\}$, and hence introduces a
second geometric scale.  The following lemma packages the uniform remainder bound needed for short-window trace
asymptotics.

\begin{lemma}[Two-scale stationary phase: uniform leading term]
\label{lem:J.twoscale}
Assume $\phi$ has a unique nondegenerate critical point $x_0$ on the set $\{\psi=0\}$, and that in a neighborhood of
$x_0$ the level set $\{\psi=0\}$ is a smooth hypersurface transversal to the gradient directions needed for the
stationary phase reduction.  Then, for $\lambda\to\infty$ and $\eta\in(0,1]$,
\begin{equation}\label{eq:J.twoscale.bound}
I(\lambda,\eta)
=
\eta\,e^{i\lambda\phi(x_0)}\lambda^{-(d-1)/2}\,A_0
+
O\!\left(\eta\,\lambda^{-(d+1)/2}\right)
+
O_N\!\left(\lambda^{-N}\eta^{-M_N}\right),
\end{equation}
where $A_0$ is an explicit constant depending on $\chi$, $\phi$, $\psi$, and $a$ at $x_0$, and $M_N$ is a finite loss
arising from derivatives of the cutoff $\chi(\psi/\eta)$.
\end{lemma}

\begin{remark}[Interpretation in localized trace regimes]
\label{rem:J.twoscale.interp}
In the localized trace setting, the admissible regime $\eta\ge \lambda^{-\theta}$ ensures that the derivative losses
$\eta^{-M_N}$ can be traded against powers of $\lambda$ by choosing $N$ large.  This is one of the mechanisms by which
the window constraint enters the error hierarchy.
\end{remark}

\section{Oscillatory integrals with endpoint singularities}
\label{sec:J.endpoints}

Several trace contributions (notably the identity term after reduction to radial variables) lead to oscillatory
integrals with mild endpoint singularities.  We record a representative estimate.

\begin{lemma}[Oscillatory integral with algebraic singularity]
\label{lem:J.endpoint}
Let $\alpha>-1$ and $A\in C^\infty([0,1])$.  Then for $\lambda\ge 1$,
\begin{equation}\label{eq:J.endpoint.bound}
\left|\int_0^1 e^{i\lambda t}\,t^\alpha\,A(t)\,dt\right|
\le
C_{\alpha,A}\,\lambda^{-\alpha-1},
\end{equation}
and if $A(0)\neq 0$ the decay rate $\lambda^{-\alpha-1}$ is sharp.
\end{lemma}

\begin{proof}
Integrate by parts once after rewriting $t^\alpha A(t)=\partial_t\big(t^{\alpha+1}B(t)\big)$ for a smooth $B$; iterate
as needed.
\end{proof}

\section{Bessel-type oscillatory transforms}
\label{sec:J.bessel}

The Kuznetsov-type applications require oscillatory integral transforms involving Bessel functions.
We record the structural estimate used in Block~8.2.

\subsection{A model Bessel transform}
\label{subsec:J.bessel.model}

Let $K_{\lambda,\eta}(x)$ denote an integral of the form
\begin{equation}\label{eq:J.bessel.transform}
K_{\lambda,\eta}(x)
=
\int_{\R} \chi_\eta(r-\lambda)\, \mathbf{J}(r;x)\,dr,
\end{equation}
where $\mathbf{J}(r;x)$ is a Bessel kernel (e.g.\ combinations of $J_{2ir}$ or $K_{2ir}$ depending on the specific
normalization of the Kuznetsov formula).

The key feature is that $\chi_\eta$ localizes $r$ to a window of width $\eta$, so that asymptotics may be obtained by
expanding $\mathbf{J}(r;x)$ around $r=\lambda$ and integrating termwise, with error controlled by bounds on
$\partial_r^k \mathbf{J}(r;x)$.

\begin{lemma}[Localized Bessel transform: schematic bound]
\label{lem:J.bessel.bound}
Assume $\mathbf{J}(r;x)$ is smooth in $r$ and satisfies for $k\le K$ the uniform derivative bound
$|\partial_r^k \mathbf{J}(r;x)|\le C_k\,\mathcal{W}(x)\,(1+|r|)^k$.
Then for $\lambda\ge 1$ and $\eta\in(0,1]$,
\begin{equation}\label{eq:J.bessel.bound}
K_{\lambda,\eta}(x)
=
\eta\,\mathbf{J}(\lambda;x)\int_{\R}\chi(u)\,du
+
O\!\left(\eta^2 \sup_{|r-\lambda|\le 2\eta} |\partial_r \mathbf{J}(r;x)|\right),
\end{equation}
and the implied constant depends only on finitely many seminorms of $\chi$.
\end{lemma}

\begin{remark}[Use in Kloosterman control]
\label{rem:J.kloosterman}
In the off-diagonal Kuznetsov analysis, the oscillation in $\mathbf{J}(r;x)$ as a function of $x$ yields additional
cancellation after summation over moduli.  Lemma~\ref{lem:J.bessel.bound} isolates the $r$-localization step; the
subsequent $x$-analysis is performed by stationary phase in the modulus variable, as in Chapter~\ref{chap:05-microlocal}.
\end{remark}

\section{A bookkeeping lemma for remainder ledgers}
\label{sec:J.ledger}

We conclude with a lemma that converts repeated uses of stationary phase into a single ledger-style bound.

\begin{lemma}[Remainder aggregation principle]
\label{lem:J.aggregation}
Let $I_k(\lambda,\eta)$, $k=1,\dots,K$, be oscillatory integrals each admitting a decomposition
\[
I_k(\lambda,\eta)=M_k(\lambda,\eta)+R_k(\lambda,\eta),
\qquad
|R_k(\lambda,\eta)|\le \mathcal{E}_k(\lambda,\eta).
\]
Then the aggregated integral $I=\sum_{k=1}^K I_k$ satisfies
\[
I(\lambda,\eta)
=
\sum_{k=1}^K M_k(\lambda,\eta)
+
R(\lambda,\eta),
\qquad
|R(\lambda,\eta)|\le \sum_{k=1}^K \mathcal{E}_k(\lambda,\eta).
\]
\end{lemma}

\begin{remark}[Interpretation]
\label{rem:J.aggregation.interp}
Lemma~\ref{lem:J.aggregation} is the formal justification for the consolidated error ledger strategy adopted in
Chapter~\ref{chap:08-normalization}: distinct sources of oscillatory remainder are tracked separately and then combined
linearly into a single quantified object.
\end{remark}

\section*{Bibliographic notes}

Stationary phase and oscillatory integral asymptotics in the form used here are standard; see, e.g., the analytic
background in semiclassical texts such as \cite{Zworski2012} and in microlocal trace formula treatments.
For Bessel transform asymptotics relevant to Kuznetsov-type formulae, see classical treatments in the automorphic
literature (e.g.\ \cite{Iwaniec2002}).

% ======================================================================
% End of J-stationary-phase-and-oscillatory-integrals.tex
% ======================================================================
